
\section{\trGenus{} $2$ 的\trCurve}

\textbf{19.6.1. 把\trGenus{} $1$ 留待日后的缘由。}\quad 此刻就直接跳到\trGenus{} $1$，或许是最顺理成章的做法，但\trEllipticCurve{}的理论格外丰富而精微，所以我们把它留到 \S19.9。

一般来说，视 $g=0$、$g=1$ 还是 $g\ge 2$ 而定，\trCurve{}会呈现出颇为不同的性态（拓扑地、算术地、几何地）。这种三分法可以延伸到更高维的\trVariety。在\trCurve{}的情形，我们对此已有了一些初步的体会。算术上，\trGenus{} $0$ 的\trCurve{}可以有许多许多\trRationalPoint，\trGenus{} $1$ 的\trCurve{}可以有许多\trRationalPoint，而由\trFaltings{}定理（\trMordell{}猜想），\trGenus{}至少为 $2$ 的任何\trCurve{}都至多只有有限多个\trRationalPoint。（于是，早在\trTaniyama--\trShimura{}猜想 [Wi, TWi] 的若干重要情形——从而\trFermat{}大定理——得到证明之前，我们就已经知道：$\mathbb{P}^2$ 中的 $x^n+y^n=z^n$ 在 $n\ge 4$ 时至多只有有限多个有理解，因为这类\trCurve{}的\trGenus{}为 $\binom{n-1}{2}>1$；见 \eqref{eq:genus-plane-curve}。）用微分几何的语言来说：\trGenus{} $0$ 的\trRiemannSurface{}具有正曲率，\trGenus{} $1$ 的\trRiemannSurface{}是平坦的，而\trGenus{}至少为 $2$ 的\trRiemannSurface{}具有负曲率。事实是：\trGenus{}至少为 $2$ 的\trCurve{}具有有限的\trAutomorphismGroup（见 \S21.4.9），而\trGenus{} $1$ 的\trCurve{}具有一维的\trAutomorphismGroup（见\mbox{问题~19.10.6}），并且\trAlgClosedField{}上唯一的那条\trGenus{} $0$ \trCurve{}具有三维的\trAutomorphismGroup（见\mbox{练习~15.5.A} 与\mbox{练习~15.5.B}）。（关于这一话题的更多内容，见\mbox{练习~21.5.F}。）

\textbf{19.6.2. 回到\trGenus{} $2$ 的\trCurve。}

在\trAlgClosedField{}上，我们在 \S19.3 中看到：\trGenus{} $0$ 的\trCurve{}只有一条。在 \S19.5 中，我们确立了：存在\trGenus{} $2$ 的超椭圆\trCurve。可我们要怎样才能把\trGenus{} $2$ 的\trCurve{}真正拿到手呢？例如，它们全都超椭圆吗？它们有「多少」条？我们现在就来着手处理这些问题。

取定一条\trGenus{}为 $g=2$ 的\trCurve{} $C$。此时 $\omega_C$ 的\trDegree{}为 $2g-2=2$，并且带有两个截面（\mbox{练习~19.2.A}）。由\mbox{练习~19.5.B}，$\omega_C$ 是\trBasePointFree{}的，从而诱导出一个二重覆盖 $C\to\mathbb{P}^1$（在 $\mathbb{P}^1$ 的\trAutomorphism{}之下唯一；这些\trAutomorphism{}我们在 \S15.5.2 中研究过）。反过来，任何一个二重覆盖 $C\to\mathbb{P}^1$ 都来自一个带有至少两个截面的、次数为 $2$ 的\trInvertibleSheaf。于是，若 $g(C)=2$，这个\trInvertibleSheaf{}就必定是\trCanonical{}层——理由无他：它是一个\trDegree{}为 $2g-2=2$、带有 $g=2$ 个截面的\trLineBundle（或者，换个思路，由\mbox{命题~19.5.6} 的最容易情形也可得出）。

于是，在\trGenus{} $2$ 的\trCurve{}与 $\mathbb{P}^1$ 的\trGenus{} $2$ 二重覆盖（在 $\mathbb{P}^1$ 的\trAutomorphism{}意义下）之间，我们便有了一个自然的\trBijection。若\trCharacteristic{}不为 $2$，则超椭圆\trRiemann--\trHurwitz{}公式（\mbox{定理~19.5.1}）表明：这个二重覆盖在 $2g+2=6$ 个几何点上分歧。特别地，我们有一「个三维的\trGenus{} $2$ \trCurve{}空间」。这并不严格，但我们当然能够证明：存在无穷多条\trGenus{} $2$ 的\trCurve。确切地说：

\textbf{19.6.A.（练习）}\quad 取定一个\trCharacteristic{}不为 $2$ 的\trAlgClosedField{} $k$。证明：$k$ 上存在无穷多条（两两）互不\trIsomorphic{}的\trGenus{} $2$ \trCurve。

\textbf{19.6.B.（练习）}\quad 证明：每条\trGenus{} $2$ 的\trCurve{}（在\trCharacteristic{}不为 $2$ 的任意域上）都具有有限的\trAutomorphismGroup。
